Gårding\textquotesingle s inequality is much easier to prove for second-order equations; i.e., in the case where $L(\bm{x},D)$ is a second-order differential operator of the form
\begin{equation}\label{eq:second-order-garding-divergence-operator}
L(\bm{x},D)u:=\sum_{i,j=1}^n\frac{\partial}{\partial x_i}a_{ij}(\bm{x})\frac{\partial u}{\partial x_j}+\sum_{i=1}^nb_i(\bm{x})\frac{\partial u}{\partial x_i}+c(\bm{x})u,
\end{equation}
with corresponding bilinear form
\begin{equation}\label{eq:second-order-garding-bilinear-form}
\begin{aligned}
B[v,u]&:=-\sum_{i,j=1}^n\int_\Omega a_{ij}(\bm{x})u_{x_j}v_{x_i}\,\mathrm{d}\bm{x}\\
&\quad+\sum_{i=1}^n\int_\Omega b_i(\bm{x})u_{x_i}v\,\mathrm{d}\bm{x}+\int_\Omega c(\bm{x})uv\,\mathrm{d}\bm{x}.
\end{aligned}
\end{equation}
In this case we do not need to use either Fourier transforms or the partition of unity technique, and the proof can be carried out under weaker hypotheses on the higher-order coefficients.

\begin{Theorem}
\label{thm:second-order-garding-inequality}
Let $\Omega$ be a bounded domain. Let $L(\bm{x},D)$ be a second-order linear partial differential operator in divergence form of the form described in \eqref{eq:second-order-garding-divergence-operator} such that for some $\theta>0$ the uniform ellipticity condition \eqref{eq:second-order-ellipticity-negative-coefficient-matrix} holds. Also suppose that $a_{ij},b_k\in L^\infty(\Omega)$ for $i,j=1,\ldots,n$, $k=0,\ldots,n$. Then there exist constants $c_3$ and $\lambda_G\geq0$ such that
\begin{equation}\label{eq:second-order-garding-energy-inequality}
B[u,u]+\lambda_G\|u\|_{L^2(\Omega)}^2\geq c_3\|u\|_{H^1(\Omega)}^2\quad\text{for all }u\in H_0^1(\Omega),
\end{equation}
where $B$ is as defined in \eqref{eq:second-order-garding-bilinear-form}.
\end{Theorem}

\begin{Proof}
We start by using the uniform ellipticity condition and Hölder\textquotesingle s inequality to get\[\begin{aligned}
B[u,u]&:=-\sum_{i,j=1}^n\int_\Omega a_{ij}(\bm{x})u_{x_i}u_{x_j}\,\mathrm{d}\bm{x}\\
&\quad+\sum_{i=1}^n\int_\Omega b_i(\bm{x})u_{x_i}u\,\mathrm{d}\bm{x}+\int_\Omega c(\bm{x})u^2\,\mathrm{d}\bm{x}\\
&\geq\theta\int_\Omega|\nabla u|^2\,\mathrm{d}\bm{x}-\max\|b_i\|_{L^\infty(\Omega)}\int_\Omega|\nabla u|\,|u|\,\mathrm{d}\bm{x}\\
&\quad-\|c\|_{L^\infty(\Omega)}\int_\Omega|u|^2\,\mathrm{d}\bm{x}.
\end{aligned}\]We now use \eqref{eq:garding-elementary-young-product-inequality} and Poincaré\textquotesingle s inequality \eqref{eq:poincare-zero-boundary-highest-derivative-control} to get\[\begin{aligned}
B[u,u]&\geq\dfrac{\theta}{2}\int_\Omega|\nabla u|^2\,\mathrm{d}\bm{x}-C\|u\|_2^2\\
&\geq C_1\|u\|_{1,2}^2-\lambda_G\|u\|_2^2.
\end{aligned}\]This completes the proof.
\end{Proof}
